\documentclass{article}
\usepackage{amsmath}
\usepackage[margin=1in]{geometry}
\usepackage{tikz}

\begin{document}
\title{ELCT 432: Problem Assignment \#6}
\author{Jacob C. Vaught}
\date{11-06-2023}
\maketitle


\subsection*{P1) Assume that \( P_{TX} = 1 \) Watt.}

\begin{itemize}
    \item[a)] What is \( P_{TX} \) in dBW?
    \item[b)] What is \( P_{TX} \) in dBm?
    \item[c)] Assume that the channel attenuates the signal by a factor \( 10^{-4} \). What is the amount of attenuation in dB?
    \item[d)] What is the received signal power in dBm?
    \item[e)] Assume that noise variance is \( 10^{-4} \). What is the noise power in dBm?
    \item[f)] What is SNR in linear scale?
    \item[g)] What is SNR in dB?
\end{itemize}

\textbf{Solution:}

\begin{itemize}
    \item[a)] Given \( P_{TX} = 1 \) W, in dBW (decibels relative to one watt), it is calculated as:
    \[ \text{dBW} = 10 \log_{10}(P_{TX}) = 10 \log_{10}(1) \]
    \[ 10 \log_{10}(1) = 10 \times 0 \]
    Therefore, \( P_{TX} = 0 \) dBW.

    \item[b)] To convert \( P_{TX} \) into dBm (decibels relative to one milliwatt), we use:
    \[ \text{dBm} = 10 \log_{10}(P_{TX} \times 1000) \]
    \[ 10 \log_{10}(1000) = 10 \times 3 \]
    Hence, \( P_{TX} = 30 \) dBm.

    \item[c)] The attenuation factor in dB is given by:
    \[ \text{Attenuation in dB} = 10 \log_{10}(10^{-4}) \]
    \[ 10 \log_{10}(10^{-4}) = 10 \times (-4) \]
    This results in \( -40 \) dB of attenuation.

    \item[d)] The received signal power in dBm after the attenuation is:
    \[ P_{RX} = P_{TX} \times 10^{-4} = 10^{-4} \text{ W} \]
    \[ \text{dBm} = 10 \log_{10}(10^{-4} \text{ W} \times 1000) \]
    \[ 10 \log_{10}(0.1) = 10 \times (-1) \]
    Therefore, \( P_{RX} = -10 \) dBm.

    \item[e)] The noise power in dBm with a noise variance of \( 10^{-4} \) is:
    \[ P_{\text{noise}} = 10^{-4} \]
    \[ P_{\text{noise in dBm}} = 10 \log_{10}(P_{\text{noise}} \times 1000) \]
    \[ 10 \log_{10}(0.1) \]
    Thus, \( P_{\text{noise}} = -10 \) dBm.

    \item[f)] The Signal-to-Noise Ratio (SNR) in linear scale is:
    \[ \text{SNR}_{\text{lin}} = \frac{P_{TX}}{P_{\text{noise}}} = -3 \]

    \item[g)] To find the SNR in dB, we calculate:
    \[ \text{SNR}_{\text{dB}} = 10 \log_{10}\left(\frac{P_{TX}}{P_{\text{noise}}}\right) \]
    \[ \text{SNR}_{\text{dB}} = 10 \log_{10}(P_{TX}) - 10 \log_{10}(P_{\text{noise}}) \]
    \[ \text{SNR}_{\text{dB}} = 30 \text{ dBm} - (-10 \text{ dBm}) \]
    Therefore, SNR in dB is \( 40 \) dB.
\end{itemize}
\newpage
\subsection*{P2}
We aim to calculate the necessary transmit power for a conventional AM system to achieve a 30 dB Signal-to-Noise Ratio (SNR) at the receiver, given a channel attenuation of 90 dB and additive white noise with a Power Spectral Density (PSD) of \( N_0 = 10^{-14} \) Watts/Hz over a bandwidth of 2 MHz. The message signal is assumed to have a power of 1 Watt, and the modulation index is 0.8.

\subsubsection*{Solution for A}
The SNR for a conventional AM system can be defined as:
\[ 10^3 = \frac{0.8^2 \times 1}{1 + 0.8^2 \times 1} \times \frac{P_t \times 10^{-9}}{2 \times 10^6 \times 10^{-14}} \]
\newline
Rearranging the formula to solve for the transmit power \( P_t \):
\[ P_t = 10^3 \times \frac{1 + 0.8^2 \times 1}{0.8^2 \times 1} \times 2 \times 10^6 \times 10^{-14} \times 10^9 \]
\newline
Solving this equation numerically, we find that the required transmit power \( P_t \) is approximately 51,250 Watts, or 47.10 dBW when converted from Watts to decibels relative to one Watt.

\subsubsection*{Solution for B}
The SNR for DSB-SC AM can be expressed as:
\[ \text{SNR}_{\text{DSB-SC}} = \frac{P_r}{P_N} \]
\noindent
Where \( P_r \) is the received power and \( P_N \) is the noise power. The received power \( P_r \) is the transmitted power \( P_t \) adjusted for channel attenuation, and the noise power \( P_N \) is the product of the noise power spectral density \( N_0 \) and the bandwidth \( B \). The equation to achieve an SNR of 30 dB (1000 in linear scale) is:
\[ \frac{P_t \times 10^{-9}}{2 \times 10^6 \times 10^{-14}} = 10^3 \]
\noindent
Solving for the transmit power \( P_t \):
\[ P_t = 10^3 \times 2 \times 10^6 \times 10^{-14} \times 10^9 \]
\noindent
Numerically, the required transmit power \( P_t \) for DSB-SC AM is found to be 20,000 Watts, which is equivalent to 43.01 dBW.


\section*{DSB-SC AM Demodulation Gain}
For DSB-SC AM, the SNR after demodulation remains unchanged from the SNR before demodulation because the process of synchronous detection does not introduce or remove noise. Therefore, the demodulation gain \( G_d \) is:
\[ G_d = \frac{\text{SNR}_{\text{after}}}{\text{SNR}_{\text{before}}} = 1 \]

\section*{Conventional AM Demodulation Gain}
In Conventional AM, the demodulation gain \( G_d \) accounts for the presence of the carrier in the received signal. The power in the sidebands is proportional to the square of the modulation index \( a \). The demodulation gain is derived as follows:

The SNR before demodulation is the ratio of the carrier power plus sideband power to the noise power:
\[ \text{SNR}_{\text{before}} = \frac{P_C(1 + a^2)}{N_0 B} \]

After demodulation, only the sidebands contribute to the message signal:
\[ \text{SNR}_{\text{after}} = \frac{a^2 P_C}{2 N_0 B} \]

Thus, the demodulation gain for Conventional AM is:
\[ G_d = \frac{\text{SNR}_{\text{after}}}{\text{SNR}_{\text{before}}} = \frac{2a^2}{1 + a^2} \]

\section*{FM Demodulation Gain}
The demodulation gain for FM is a result of the capture effect and the noise-quenching properties inherent to FM. The derivation of the demodulation gain \( G_d \) for FM involves the following steps:

The pre-detection SNR is the ratio of the carrier power to the noise power:
\[ \text{SNR}_{\text{before}} = \frac{P_C}{N_0 B_c} \]
where \( B_c \) is the bandwidth of the FM signal, typically approximated by Carson's rule.

After demodulation, the SNR is improved due to the constant envelope nature of FM and its resistance to amplitude noise. The improvement is proportional to the square of the modulation index \( \beta \), which is the ratio of the frequency deviation to the modulating frequency:
\[ \text{SNR}_{\text{after}} = \text{SNR}_{\text{before}} \times \beta^2 \times \frac{f_m}{N_0} \]

The modulation index for FM is given by:
\[ \beta = \frac{\Delta f}{f_m} \]

Combining these relations, the demodulation gain for FM is found to be:
\[ G_d = \frac{\text{SNR}_{\text{after}}}{\text{SNR}_{\text{before}}} = \beta^2 \times \frac{f_m}{N_0} \]
which, under the assumption of a high modulation index and a wideband FM system, simplifies to:
\[ G_d = 3\beta^2(\beta + 1) \]

\end{document}
